\chapter{Développement d'une fonctionnalité en code de production}
\label{chap:t7}
% Entreprise : Beemo Technologie — Volet « développer dans un existant endetté »

\section{Contexte et besoin}
% TODO : présenter la feature demandée, le besoin métier/utilisateur, et
% son insertion dans la base de code existante (legacy ? bien outillée ?).
Le métier de Beemo n'est pas de proposer un \textbf{service de cloud} à la
manière des grands fournisseurs — Amazon, Google ou Microsoft. Il s'agit de
fournir aux \textbf{professionnels} — hébergeurs et entreprises pour qui la
donnée constitue un actif critique — un moyen de \textbf{sécuriser leurs
données} et de garantir leur \textbf{restauration} en cas d'incident, qu'il soit
d'origine \textbf{physique} ou \textbf{malveillante}. Le produit central, la
\textbf{BBox}, est une appliance de \textbf{sauvegarde} qui peut être
\textbf{externalisée} hors site afin de survivre à la destruction du système
d'origine.

L'un des scénarios de perte de données les plus redoutés est le
\textbf{ransomware}, qui chiffre les données d'un système pour en exiger une
rançon. La BBox répond directement à cette menace en conservant des
\textbf{jeux de sauvegarde} hors d'atteinte du système compromis. Mais une
question demeure : \textbf{que se passe-t-il si la BBox elle-même est
compromise} ? Un attaquant parvenu jusqu'à l'appliance pourrait alors tenter
d'\textbf{altérer les jeux de sauvegarde} ou de \textbf{saboter leur
restauration} — anéantissant la dernière ligne de défense du client.

C'est pour répondre à ce scénario que Beemo développe une nouvelle
fonctionnalité, l'\textbf{immuabilité}. Son principe : rendre un jeu de
sauvegarde \textbf{insupprimable} et \textbf{inaltérable} tant qu'un
\textbf{délai de rétention} n'est pas écoulé, y compris pour un utilisateur
disposant de droits sur la box — et donc pour un attaquant qui en aurait pris le
contrôle. J'ai participé à son développement aux côtés de \textbf{deux autres
développeurs}, en me concentrant sur la partie \textbf{embarquée sur la box},
écrite en \textbf{C++}.

\section{Objectif de la mission}
% TODO : implémenter la fonctionnalité en code de production, en limitant
% l'ajout de dette technique (qualité, tests, intégration propre).
L'objectif n'était pas seulement de \textbf{livrer la fonctionnalité}, mais de
le faire dans un contexte particulièrement défavorable : celui d'une base de
code embarquée portant une \textbf{dette technique considérable}. L'intérêt de
cette tâche, au regard du fil rouge de ce mémoire, tient précisément là —
\textbf{comment injecter une fonctionnalité moderne dans du code dégradé sans
en briser l'existant} ?

Le firmware de la box cumule en effet la plupart des symptômes de la dette :
\begin{itemize}
  \item un \textbf{abus de directives de préprocesseur} (\texttt{\#ifdef},
        \texttt{\#ifndef}) pour gérer les multiples cas de \textbf{compilation}
        de la box, rendant le code difficile à lire et à raisonner ;
  \item des \textbf{fonctions extrêmement longues} et une \textbf{maîtrise
        approximative du C++}, donnant lieu à des constructions trompeuses —
        par exemple des contrôles du type \texttt{this == nullptr}, un cas qui
        \textbf{ne peut jamais se produire} et trahit une incompréhension du
        langage ;
  \item une \textbf{absence totale de tests} — ni \gls{test-regression}, ni
        tests unitaires — qui prive de tout filet de sécurité lors d'une
        modification ;
  \item un fichier \texttt{main.h} qui \textbf{inclut l'ensemble des fichiers}
        du projet, créant un \textbf{couplage massif} : réutiliser une
        fonctionnalité existante depuis du code neuf devient une opération
        complexe et risquée.
\end{itemize}
La contrainte centrale était donc de \textbf{ne pas aggraver cette dette} : il
fallait écrire du code \textbf{moderne, testable et isolé}, tout en
\textbf{dialoguant} avec un existant qui ne s'y prête pas.

\section{Conception et intégration}
% TODO : choix de conception, impact sur l'architecture existante,
% interaction avec le travail des tâches précédentes (tests, build).
La réponse à cette tension repose sur une \textbf{couche intermédiaire}
d'\textbf{adaptation}, conçue par David, l'un des développeurs de l'équipe.
L'idée est d'\textbf{encapsuler} une fonctionnalité — qu'elle soit
\textbf{legacy} ou \textbf{nouvelle} — derrière une \textbf{classe} pouvant être
appelée \textbf{indifféremment depuis l'ancien code ou depuis le code neuf}. Les
appels vers l'existant sont \textbf{masqués} par des \textbf{pointeurs} transmis
à cette classe, ce qui \textbf{découple} le nouveau code du couplage imposé par
\texttt{main.h}.

Ce mécanisme d'\emph{enrobage} (\emph{wrapping}) joue le rôle de
\textbf{frontière} : d'un côté, le firmware historique et ses contraintes ; de
l'autre, un espace où l'on peut écrire du \textbf{C++ réellement moderne} pour
réaliser la fonctionnalité. La fonctionnalité d'immuabilité a ainsi pu être
bâtie autour de deux briques propres :
\begin{itemize}
  \item un \textbf{automate à états} (\emph{state machine}) qui pilote le cycle
        de vie d'un jeu de sauvegarde immuable — création, période de rétention,
        mise en pause du \textbf{planificateur de sauvegarde}, expiration ;
  \item un \textbf{service de chiffrement} qui garantit l'\textbf{intégrité} des
        données. La base est \textbf{lue et déchiffrée à l'aide d'une clé} : si
        elle a été \textbf{altérée}, le déchiffrement échoue, \textbf{le code le
        détecte} et la box se \textbf{bloque} — empêchant toute manipulation des
        jeux de données réels par un attaquant parvenu jusqu'à l'appliance.
\end{itemize}
L'immuabilité vient ainsi s'ajouter aux mécanismes de sécurisation existants :
elle \textbf{empêche la suppression} d'un jeu de sauvegarde considéré comme
immuable tant que le \textbf{délai de rétention} n'a pas été respecté.

\section{Réalisation}
% TODO : extrait de code de production commenté + éventuels tests associés.
La couche d'adaptation est le point d'entrée de toute la fonctionnalité.
L'extrait ci-dessous en illustre le principe : une classe moderne expose le
service d'immuabilité au reste du système, tout en dissimulant les appels au
firmware historique derrière un pointeur opaque transmis à la construction.

\begin{lstlisting}[language=C++,caption={Couche d'adaptation isolant le code neuf du firmware legacy : les appels à l'existant sont masques derriere un pointeur (extrait illustratif).},label={lst:t7}]
// La fonctionnalite d'immuabilite est ecrite en C++ moderne, isolee du
// firmware historique. Le couplage impose par main.h est neutralise :
// le code neuf ne connait l'existant qu'a travers un pointeur opaque.
class ImmutabilityService {
public:
    // Le legacy est injecte par pointeur a la construction : le code neuf
    // n'inclut jamais main.h et reste testable de maniere isolee.
    explicit ImmutabilityService(LegacyBackupContext* ctx) noexcept
        : m_ctx(ctx) {}

    // Verrouille un jeu de sauvegarde : il devient insupprimable et
    // inalterable tant que le delai de retention n'est pas ecoule.
    Result lock(const BackupSetId& id, RetentionDelay delay) {
        // L'automate a etats pilote le cycle de vie de l'immuabilite.
        if (!m_state.canTransitionTo(State::Locked))
            return Result::Rejected;

        // L'integrite est garantie par chiffrement : a la relecture, un
        // echec de dechiffrement revele toute alteration de la base.
        auto sealed = m_cipher.seal(id, delay);
        m_state.transition(State::Locked);
        return sealed ? Result::Ok : Result::IntegrityError;
    }

private:
    LegacyBackupContext* m_ctx;   // pointeur opaque vers le firmware legacy
    ImmutabilityStateMachine m_state;
    CipherService m_cipher;
};
\end{lstlisting}

\section{Difficultés rencontrées}
% TODO : contraintes du code existant, revues, exigences de qualité.
La principale difficulté n'a pas été la \textbf{conception de la fonctionnalité}
en elle-même, mais sa \textbf{cohabitation avec l'existant}. Comprendre le
firmware historique — afin de \textbf{réutiliser ses fonctionnalités} depuis le
code neuf — s'est avéré coûteux, en raison du \textbf{couplage par
\texttt{main.h}}, des \textbf{fonctions interminables} et des \textbf{cas de
compilation} multipliés par les directives de préprocesseur.

L'\textbf{absence de tests} a constitué un risque permanent : sans
\gls{test-regression}, la moindre modification dans une zone partagée pouvait
\textbf{casser silencieusement} une fonctionnalité existante. La couche
d'adaptation a ici joué un rôle de \textbf{garde-fou}, en \textbf{circonscrivant}
les points de contact avec le legacy et en rendant le code neuf
\textbf{testable de manière isolée} — là où le reste du firmware ne l'est pas.

\section{Résultats et bénéfice business}
% TODO : valeur livrée, qualité maintenue (pas de dette ajoutée),
% éventuelle couverture de tests. Relier au fil rouge.
Cette tâche illustre une facette essentielle de la lutte contre la dette
technique : il ne s'agit pas toujours de \textbf{résorber} l'existant, mais
parfois de savoir \textbf{construire proprement à côté} sans l'aggraver. Grâce à
la couche d'adaptation, la fonctionnalité d'immuabilité a pu être développée en
\textbf{C++ moderne}, \textbf{isolée} et \textbf{testable}, alors même qu'elle
s'insère dans un firmware fortement dégradé et dépourvu de filet de sécurité.

Sur le plan \textbf{business}, l'immuabilité renforce directement la
\textbf{proposition de valeur} de Beemo : elle protège les jeux de sauvegarde
non seulement contre la compromission du système client, mais aussi contre la
\textbf{compromission de la BBox elle-même}. En garantissant qu'un jeu de
sauvegarde immuable ne peut être ni supprimé ni altéré avant l'expiration de son
\textbf{délai de rétention}, et en \textbf{bloquant} la box au moindre signe
d'altération détecté par le chiffrement, la fonctionnalité ferme l'un des
derniers angles d'attaque possibles — consolidant la promesse fondamentale du
produit : \textbf{la donnée sauvegardée n'est jamais perdue}.
